
% ════════════════════════════════════════════
% 问题重述（1.问题重述.tex）
% ════════════════════════════════════════════
\section{问题重述}

\subsection{问题背景}

光伏出力受天气影响，具有间歇性与随机性，大规模接入会给电网功率平衡带来压力。微网聚合分布式电源、储能与本地负荷\cite{luzongxiang2007}，通过本地消纳光伏、平抑负荷波动减少对公共电网的依赖；非孤岛运行时，微网既可外购电力，也可将余电反馈电网或存入储能\cite{yangxinfa2014,wuxiong2014}，其能量管理与运行调度是微网研究的重点。

光伏只在白天出力且峰谷差大，小区负荷与光伏出力在一天、一年内均随时间变化；外网电价存在峰谷价差，波动电价下价差更大\cite{zhangqin2008}。若按负荷实时外购，峰时段费用很高；引入储能后，可在低价或光伏有剩余时充电、高价或光伏不足时放电，降低购电费用\cite{zhoulin2011,sangbingyu2014}。但储能容量、功率与效率有限，购电计划又必须在当天负荷与光伏实际发生之前制定。

\subsection{问题重述}

题目给出一座含光伏与储能的小区微网，其储能容量为 12000 kWh、最大充放电功率 5000 kW、充放电效率 90\%，储电量须保持在 1200$\sim$10800 kWh，2025 年 1 月 1 日 0:00 初始储电量为 6000 kWh；负荷、光伏与电价数据均以 10 分钟为间隔给出。据此建立数学模型解决以下四个问题：

\begin{enumerate}
  \item 在当日负荷、光伏出力与分时电价已知的条件下，制定全天各时段的购电与充放电计划，使单日购电费用最小，并给出储能参数的可用区间；
  \item 实际运行前只能获得预测信息，需按“预测—决策—结算”的时序安排全年逐日的购电与充放电，使全年购电支出尽可能低，并对预测失准造成的供电缺口支付 5 倍电价的紧急购电费用；
  \item 光伏预报可在 0:00、6:00、12:00、18:00 四次更新，需在预报更新时滚动调整购电计划，并在计划调整代价与购电费用之间权衡；
  \item 将分时电价替换为逐时段波动的实时电价，重新评估上述策略，比较固定电价与波动电价下的费用差异。
\end{enumerate}
